\begin{figure}[htbp]
\centering
\begin{tikzpicture}[x=1cm,y=1cm,font=\small]
\foreach \x/\which/\title/\n/\pairs in {0/0/完全重叠/2/4,5/1/部分重叠/3/6,10/2/完全不重叠/4/8}{
 \begin{scope}[xshift=\x cm]
 \node at (2.2,3.25) {\title};
 \foreach \j/\name in {0/A,1/B,2/C,3/D}{\node[font=\footnotesize] at (.92+\j*.82,2.73) {\name};}
 \foreach \i in {0,1}{
  \node[anchor=east,font=\footnotesize] at (.4,2.16-\i*.65) {$q_\i$};
  \foreach \j in {0,1,2,3}{
   \draw[gridline,fill=white] (.55+\j*.82,1.9-\i*.65) rectangle ++(.74,.52);
   \ifnum\j<\n
    \draw[gridline,fill=waste] (.55+\j*.82,1.9-\i*.65) rectangle ++(.74,.52);
    \draw[gridline] (.62+\j*.82,1.97-\i*.65)--++(.60,.38);
    \draw[gridline] (.62+\j*.82,2.35-\i*.65)--++(.60,-.38);
   \fi
  }
 }
 \foreach \j in {0,1}{\filldraw[accent,fill=selected] (.55+\j*.82,1.9) rectangle ++(.74,.52);}
 \ifnum\which=0
  \foreach \j in {0,1}{\filldraw[accent,fill=selected] (.55+\j*.82,1.25) rectangle ++(.74,.52);}
 \else\ifnum\which=1
  \foreach \j in {1,2}{\filldraw[accent,fill=selected] (.55+\j*.82,1.25) rectangle ++(.74,.52);}
 \else
  \foreach \j in {2,3}{\filldraw[accent,fill=selected] (.55+\j*.82,1.25) rectangle ++(.74,.52);}
 \fi\fi
 \node[font=\footnotesize,align=center] at (2.15,.57) {读取块数：$4\to\n$\\计算格子：$4\to\pairs$};
 \end{scope}
}
\node[font=\footnotesize,text=muted] at (7.2,-.28) {蓝色：有效计算\qquad 灰色带叉：被 mask 的计算\qquad 白色：不计算};
\end{tikzpicture}
\caption{两个 query 各选两块。分组后，读取的块数和计算量随重叠程度变化。}
\label{fig:union-overlap}
\end{figure}
